\section{Explicit Smith kernel coordinates}

Fix a degree $n$ and compute a strong Smith result
\[
  Ud_nV=S.
\]
Let $r$ be the length of the nonzero diagonal prefix of $S$.  The strong
certificate gives both $V^{-1}V=I$ and $VV^{-1}=I$, so no determinant or
choice of a later inverse is required.

The Smith predicate proves that the first $r$ diagonal entries are nonzero
and that every off-diagonal entry vanishes.  A vector annihilated by $S$
therefore has zero first $r$ coordinates.  Conversely, every vector with
that zero prefix is annihilated by $S$.  Multiplication by $V$ transports
the remaining $\rank C_n-r$ coordinates to $\ker d_n$.

Lean packages this construction as
\[
  \code{kernelBasisEquiv}\ n :
  \Z^{\rank C_n-r}\simeq_{\Z}\ker d_n.
\]
Its forward map extends kernel coordinates by a zero prefix and multiplies
by $V$.  Its inverse multiplies a cycle by $V^{-1}$ and restricts to the
tail.  The two inverse proofs use the two stored inverse identities.

To express the incoming boundary in this basis, define
\[
  C=V^{-1}d_{n+1}.
\]
The transformation equation and the chain condition give
\[
  SC
  =Ud_nVV^{-1}d_{n+1}
  =U(d_nd_{n+1})=0.
\]
The nonzero Smith prefix then forces the first $r$ rows of $C$ to vanish.
The remaining rows form \code{kernelCoordinateMatrix}.  The theorem
\code{kernelBasisEquiv\_symm\_boundary} states pointwise that applying the
inverse kernel-basis equivalence to an incoming boundary is exactly
multiplication by this restricted matrix.  This is the central bridge from
matrix computation to the intrinsic quotient.
